% ============================================================
%  Conclusion Générale
% ============================================================
\chapter*{Conclusion Générale}
\addcontentsline{toc}{chapter}{Conclusion Générale}
\markboth{Conclusion Générale}{}

Ce stage a été consacré à la conception, à la réalisation et au déploiement d'un
système RAG (\eng{Retrieval-Augmented Generation}) permettant d'interroger en
langage naturel un corpus de documents d'entreprise, avec des réponses sourcées
et vérifiables. Le travail a été conduit selon la méthodologie CRISP-DM, dont
chaque phase a été traitée successivement : compréhension du besoin métier,
exploration et préparation d'un corpus de 2\,676 documents, construction de la
chaîne de traitement, évaluation de sa fiabilité et déploiement conteneurisé.

\section*{Bilan du travail réalisé}

Sur le plan fonctionnel, la chaîne complète est opérationnelle. Le pipeline
d'ingestion traite les cinq formats documentaires courants, extrait le texte,
l'enrichit de métadonnées de traçabilité et le découpe en segments : les 2\,676
fichiers du corpus ont été intégralement traités, sans aucun échec d'extraction,
produisant 2\,929 segments vectorisés et persistés dans ChromaDB. La recherche
sémantique s'exécute en 182~millisecondes en moyenne. La génération est assurée
localement par Mistral 7B, contraint par une invite qui lui interdit de sortir du
contexte fourni. L'ensemble est exposé par une API REST à quatre
\eng{endpoints} et une interface conversationnelle, déployés en trois conteneurs
Docker.

La contrainte structurante du projet a été respectée sans concession :
\textbf{aucune donnée ne quitte l'infrastructure}. Vectorisation, indexation,
recherche, génération et même l'évaluation par modèle juge s'exécutent
localement, sans clé d'API ni coût récurrent. Cette propriété, qui aurait pu
paraître une contrainte appauvrissante, s'est révélée le principal argument de
l'architecture retenue face aux solutions du marché.

Sur le plan de la traçabilité, l'objectif est atteint. Chaque réponse expose les
fichiers sources mobilisés, l'identifiant précis de chaque segment et son score
de pertinence. L'utilisateur peut vérifier ce que le système avance plutôt que de
devoir lui faire confiance --- ce qui constitue, en contexte professionnel, la
condition d'adoption d'un outil de ce type.

\section*{Principal apport de la démarche d'évaluation}

Le résultat le plus instructif du stage n'est pas la mise en service du système,
mais ce que sa mesure a révélé. À l'usage, l'assistant produisait des réponses
d'apparence satisfaisante. L'évaluation quantitative de la composante de
recherche a pourtant montré que les scores de pertinence étaient anormalement
faibles sur des questions posées en français, et que le système s'appuyait
fréquemment sur des documents inadaptés.

Une expérience contrôlée a permis d'en identifier et d'en quantifier la cause :
à contenu sémantique identique, la traduction d'une question en anglais
\textbf{triple le score de pertinence} (0,461 contre 0,153 en moyenne). Le modèle
\texttt{all-MiniLM-L6-v2}, entraîné essentiellement sur de l'anglais, projette
mal les questions françaises dans l'espace sémantique de documents anglophones.

Ce constat porte un enseignement méthodologique qui dépasse le cas particulier :
\textbf{un système RAG peut paraître fonctionner tout en s'appuyant sur les
mauvais documents}. Une réponse fluide, bien structurée et accompagnée de sources
n'est pas pour autant une réponse correcte. Seule la mesure permet de distinguer
les deux, et c'est ce qui justifie la place centrale que CRISP-DM accorde à la
phase d'évaluation avant tout déploiement.

\section*{Limites du travail}

Trois limites doivent être clairement énoncées.

La \textbf{campagne d'évaluation RAGAS n'a pas été menée à son terme}. Le module
de mesure est développé et configuré pour fonctionner avec un modèle juge local,
mais son exécution requiert la constitution préalable d'un jeu de questions
annotées de réponses de référence rédigées manuellement. Les quatre métriques ---
\eng{faithfulness}, \eng{answer relevancy}, \eng{context precision} et
\eng{context recall} --- restent donc à mesurer. Le choix a été fait de laisser
ces valeurs explicitement vacantes dans le présent rapport plutôt que de les
estimer.

Le \textbf{corpus est exclusivement transactionnel}. Les 2\,676 documents sont
des factures, bons de commande, commandes et états de stock, c'est-à-dire des
enregistrements structurés. Ils permettent de valider techniquement le pipeline,
mais ils sous-exploitent l'intérêt propre d'une architecture RAG, qui donne son
plein rendement sur des documents rédigés où l'information est diffuse. Sur des
enregistrements structurés, une requête SQL serait plus exacte et plus rapide.
Cette composition explique par ailleurs le caractère quasi inopérant du paramètre
de découpage, la longueur moyenne d'un document (498 caractères) étant très
inférieure à la taille de segment configurée (1\,000 caractères).

La \textbf{latence de génération est excessive} pour un usage quotidien : environ
110 secondes par réponse. L'inférence d'un modèle de 7 milliards de paramètres
sur CPU représente plus de 99~\% du temps de réponse total, la recherche
vectorielle étant quant à elle négligeable.

\section*{Perspectives}

Les travaux futurs se hiérarchisent naturellement à partir des mesures obtenues.

\paragraph{Adopter un modèle d'\eng{embeddings} multilingue.}
C'est l'action de rang le plus élevé, car la plus étayée expérimentalement. Le
remplacement de \texttt{all-MiniLM-L6-v2} par
\texttt{paraphrase-multilingual-MiniLM-L12-v2}, entraîné sur plus de cinquante
langues, devrait résoudre l'asymétrie linguistique documentée au
chapitre~\ref{ch:evaluation}. L'opération impose une ré-indexation complète du
corpus.

\paragraph{Conduire la campagne RAGAS.}
Une fois le modèle multilingue en place, la constitution d'un jeu de vingt
questions annotées permettra d'établir une mesure de référence des quatre
métriques, et surtout de comparer objectivement les configurations successives
plutôt que de se fier à une impression d'usage.

\paragraph{Mettre en place une recherche hybride.}
La combinaison d'une composante lexicale exacte (BM25) et de la composante
vectorielle, avec fusion des classements, lèverait la limite documentée sur la
recherche par identifiant exact --- besoin fréquent sur un corpus transactionnel
où l'utilisateur connaît souvent le numéro de la pièce qu'il cherche.

\paragraph{Ajouter un ré-ordonnancement.}
L'insertion d'un \eng{cross-encoder} après la recherche vectorielle, chargé de
réévaluer finement la pertinence des extraits candidats, améliorerait la
\eng{context precision} et permettrait d'établir enfin un seuil de rejet fiable.

\paragraph{Enrichir le corpus de documents rédigés.}
L'ajout de procédures, manuels et notes d'organisation permettrait d'exercer le
système sur le type de questions pour lequel une architecture RAG est réellement
supérieure aux alternatives, et redonnerait au réglage du découpage son rôle de
levier d'optimisation.

\paragraph{Réduire la latence.}
Le passage sur GPU, l'emploi d'un modèle plus fortement quantifié ou plus petit,
et la diffusion progressive de la réponse constituent trois leviers cumulables.

\vspace{0.5cm}

En définitive, ce stage a permis de mener un projet de science des données de la
formulation du besoin jusqu'à un système déployé et mesurable, en abordant
l'ensemble de la chaîne : ingénierie documentaire, recherche sémantique,
ingénierie d'invite, évaluation de systèmes génératifs et déploiement
conteneurisé. Au-delà des compétences techniques acquises, l'apport principal
réside dans une conviction méthodologique : dans un système fondé sur un modèle
génératif, la qualité perçue est un indicateur trompeur, et l'instrumentation de
la mesure n'est pas une étape finale de validation mais la condition même du
progrès.
